\clearpage
\section{Аналитическая часть}

\subsection{Постановка задачи}

Даны натуральные числа $k$ и $s$. Требуется определить, сколько существует
$k$-значных натуральных чисел, сумма цифр которых равна $s$. Запись натурального
числа не может начинаться с цифры 0. Например, при $k=3$ и $s=15$ ответ равен 69.
Если $s>9k$, искомых чисел нет, поскольку каждая цифра не превосходит 9.

Ответ не превосходит количества всех $k$-значных чисел $9\cdot10^{k-1}$.
При $k\le19$ это значение меньше $2^{63}$, поэтому ответ представим
64-битовым целым числом со знаком; для б\'{о}льших $k$ такое представление
может переполниться.

Нерекурсивный алгоритм решает задачу методом динамического
программирования: ответ составляется
из ответов на подзадачи меньшего размера, каждая подзадача решается один раз,
а её результат сохраняется для повторного использования~\cite{laaksonen}.
Обозначим через $F(i,j)$ количество $i$-значных натуральных чисел, сумма цифр
которых равна~$j$. Каждое $i$-значное число при $i\ge2$ получается приписыванием
справа цифры $c$ к $(i-1)$-значному числу с суммой цифр $j-c$, откуда
\begin{equation}
    \label{eq:f}
    F(i,j)=
    \begin{cases}
        1, & i=0,\ j=0,\\
        0, & i=0,\ j>0,\\
        0, & i\ge1,\ j=0,\\
        \sum\limits_{c=0}^{\min(9,j)}F(i-1,j-c), & i\ge1,\ j\ge1,
    \end{cases}
\end{equation}
где $c$~--- последняя цифра числа. Значение $F(0,0)=1$ соответствует пустой
записи. При $i=1$ в сумме отлична от нуля только величина $F(0,0)$, поэтому
$F(1,j)=1$ для $1\le j\le9$, а первая цифра числа всегда отлична от нуля.
Ответ задачи~--- $F(k,s)$.

На рекурсивный алгоритм наложены ограничения задания: нельзя использовать циклы,
массивы, строки и срезы. Поэтому сумма в формуле~\eqref{eq:f} заменяется
рекурсивной функцией
\begin{equation}
    \label{eq:s}
    S(i,j,c)=
    \begin{cases}
        0, & c>9 \text{ или } c>j,\\
        F(i-1,j-c)+S(i,j,c+1), & \text{иначе},
    \end{cases}
\end{equation}
где $S(i,j,c)$~--- сумма величин $F(i-1,j-d)$ для цифр $d$ от $c$ до $\min(9,j)$.
Тогда при $i\ge1$ и $j\ge1$ выполняется $F(i,j)=S(i,j,0)$.

\subsection{Графовые модели алгоритмов}

Алгоритм рассматривается как совокупность операторов. Однократное выполнение
оператора называется его срабатыванием. Между операторами алгоритма определяются
два отношения~\cite{voevodin}:
\begin{itemize}
    \item операционное отношение: оператор $v$ связан с оператором $u$, если $v$
    может быть выполнен непосредственно после $u$, то есть управление
    передаётся от $u$ к $v$;
    \item информационное отношение: оператор $v$ связан с оператором $u$, если $v$
    использует в качестве аргумента величину, вычисленную $u$, и эта величина
    не изменяется между их срабатываниями.
\end{itemize}

На основе этих отношений строятся четыре графовые модели~\cite{voevodin}:
\begin{itemize}
    \item граф управления~--- вершины соответствуют операторам, дуги~---
    операционному отношению;
    \item информационный граф~--- вершины соответствуют операторам, дуги~---
    информационному отношению;
    \item операционная история~--- вершины соответствуют срабатываниям операторов
    при конкретных входных данных, дуги указывают порядок срабатываний;
    \item информационная история~--- вершины соответствуют срабатываниям
    операторов при конкретных входных данных, дуги~--- информационному отношению
    между срабатываниями.
\end{itemize}

Граф управления и информационный граф описывают все возможные связи
операторов, а истории~--- одно выполнение алгоритма. Информационная история
не содержит циклов. Ярусом её вершины называется число вершин на наибольшем
пути, оканчивающемся в этой вершине. Срабатывания одного яруса не связаны
информационным отношением и могут выполняться параллельно, а число ярусов
равно наименьшему числу последовательных шагов при неограниченном числе
исполнителей и одинаковой длительности срабатываний.
